\scanpage{132}
% archive sheet number: 131

2) \deleted{On voit que} \inserted{Posons}

\[
L=H/k\inserted{\simeq B/H'},\quad L'=H'/k\simeq B/H
\]

\deleted{et considérons une forme}

\[
\deleted{\lambda:L\otimes L'\unsure{-|\longrightarrow} k}
\]

\deleted{Si $\Delta_B\in\underline{Q}_B$, on peut définir}

On a

\[
\begin{aligned}
\deleted{\unsure{\Delta_B\to B}}\ \deleted{\unsure{B\to L\otimes{}}}\quad
(L\otimes L')^\vee&\xrightarrow{\sim}\text{Esp. des formes}\\
&\quad\text{bil. \inserted{sym} sur }B\text{ nulles}\\
&\quad\text{sur }H'\deleted{\unsure{B}}\text{ et }H\deleted{\unsure{B}}
\end{aligned}
\]

(i.e. qui passent au
quotient dans $B/k\simeq L\oplus L'$ \inserted{car $B=H+H'$}
et qui sont nulles
sur $L$, et sur $L'$
donc de la forme

\[
\begin{aligned}
(\underset{\rotatebox{90}{$\in$}\,L}{x},\underset{\rotatebox{90}{$\in$}\,L'}{y}\deleted{\unsure{'}}),\
\deleted{(\unsure{x',y'})}\inserted{(\xi,\unsure{\eta|y'})}
 &\longrightarrow\lambda(x\otimes\unsure{\eta|y})\\
 &\quad+\lambda(\xi\otimes y)
\end{aligned}
\]

avec $\lambda:L\otimes L'\longrightarrow k$

Donc \deleted{$\Delta$}$\unsure{\underline{\Phi}_B|\Phi_B}$ est un torseur sous $(L\otimes L')^\vee$, et de
même pour $\underline{\Psi}_B$. Itou pour $\Delta_B$, car

\[
\begin{aligned}
(L\otimes L')^\vee&\simeq\text{formes quadratiques sur }B\\
&\qquad\text{qui sont nulles sur }H,H'
\end{aligned}
\]

(donc passent au quotient
dans $B/k\simeq L\oplus L'$ en des
formes quadratiques nulles
sur $L,L'$ \deleted{\unsure{donc formes de|la forme}} donc de
la forme

\[
\Delta(x\oplus y)=\lambda(x\otimes y)
\]

Les applications $\beta,\gamma$ sont compatibles
avec ces structures de torseurs. Cela
suffit à assurer qu'elles sont iso.

